% !TeX program = xelatex
% 47 pages: 38 presentation slides + 7-part contents + 7 part dividers + dedicated closing
\documentclass[aspectratio=169,10pt]{beamer}
\usepackage{fontspec}
\usepackage[vietnamese]{babel}
\usepackage{amsmath,amssymb,mathtools}
\usepackage{tikz}
\usetikzlibrary{arrows.meta,positioning,calc,shapes.geometric}
\usepackage{booktabs,tabularx,array,colortbl,multirow}
\usepackage{graphicx}
\usepackage{microtype}
\setmainfont{Noto Serif}
\setsansfont{Noto Sans}
\setmonofont{DejaVu Sans Mono}
\usefonttheme{serif}
\linespread{1.0}

% ---------- Red academic theme (according to user's specification) ----------
\definecolor{mainred}{HTML}{B71922}
\definecolor{darkred}{HTML}{82151B}
\definecolor{footerred}{HTML}{991820}
\definecolor{softred}{HTML}{F9E9EA}
\definecolor{darkgray}{HTML}{242424}
\definecolor{midgray}{HTML}{666666}
\definecolor{lightgray}{HTML}{F2F2F0}
\definecolor{paperbg}{HTML}{FAFAF8}
\definecolor{goodgreen}{HTML}{2E7D32}
\definecolor{warnorange}{HTML}{B85C00}
\definecolor{badred}{HTML}{7F1016}

\setbeamercolor{background canvas}{bg=paperbg}
\setbeamercolor{normal text}{fg=darkgray,bg=paperbg}
\setbeamercolor{structure}{fg=mainred}
\setbeamercolor{frametitle}{fg=white,bg=mainred}
\setbeamercolor{title}{fg=mainred}
\setbeamercolor{subtitle}{fg=midgray}
\setbeamercolor{block title}{fg=white,bg=mainred}
\setbeamercolor{block body}{fg=darkgray,bg=softred}
\setbeamercolor{alerted text}{fg=badred}
\setbeamercolor{author in head/foot}{bg=darkred,fg=white}
\setbeamercolor{title in head/foot}{bg=footerred,fg=white}
\setbeamercolor{date in head/foot}{bg=mainred,fg=white}
\setbeamerfont{frametitle}{size=\large,series=\bfseries,family=\rmfamily}
\setbeamerfont{title}{size=\LARGE,series=\bfseries,family=\rmfamily}
\setbeamerfont{block title}{series=\bfseries,family=\rmfamily}
\setbeamerfont{normal text}{family=\rmfamily}
\setbeamertemplate{navigation symbols}{}
\setbeamertemplate{headline}{}
\setbeamertemplate{itemize item}{\color{mainred}\small$\bullet$}
\setbeamertemplate{itemize subitem}{\color{midgray}\scriptsize$\circ$}
\setbeamertemplate{blocks}[rounded][shadow=false]
\setbeamersize{text margin left=0.72cm,text margin right=0.72cm}
\setlength{\parskip}{0.0em}
% Top academic section breadcrumb + presentation title.
\setbeamercolor{sectionheader}{fg=darkred,bg=softred}
\setbeamertemplate{frametitle}{%
  \nointerlineskip%
  \begin{beamercolorbox}[wd=\paperwidth,ht=0.23cm,dp=0.08cm,leftskip=0.69cm,rightskip=0.68cm]{sectionheader}%
    {\fontsize{7.4}{8.6}\selectfont\bfseries\insertsectionhead\hfill CHUYÊN ĐỀ 2 \enspace / \enspace USD/VND}%
  \end{beamercolorbox}%
  \nointerlineskip%
  \begin{beamercolorbox}[wd=\paperwidth,ht=0.66cm,dp=0.17cm,leftskip=0.68cm]{frametitle}%
    \usebeamerfont{frametitle}\insertframetitle
  \end{beamercolorbox}%
}
% References live in a dedicated right-aligned baseline ABOVE the footer.
% This placement does not consume space in the main slide flow.
\newcommand{\slideRef}{}
\AtBeginEnvironment{frame}{\gdef\slideRef{}}
\setbeamercolor{reference bar}{bg=paperbg,fg=midgray}
\setbeamertemplate{footline}{%
  \leavevmode%
  \vbox{%
  \hbox{\begin{beamercolorbox}[wd=\paperwidth,ht=1.65ex,dp=0.48ex,rightskip=0.66cm]{reference bar}%
    \hfill\fontsize{6.5}{7.5}\selectfont\color{midgray}\slideRef
  \end{beamercolorbox}}%
  \hbox{%
  \begin{beamercolorbox}[wd=.23\paperwidth,ht=2.65ex,dp=1.05ex,leftskip=0.35cm]{author in head/foot}%
    \color{white}\scriptsize\bfseries CHUYÊN ĐỀ 2
  \end{beamercolorbox}%
  \begin{beamercolorbox}[wd=.64\paperwidth,ht=2.65ex,dp=1.05ex,center]{title in head/foot}%
    \color{white}\scriptsize PHÂN TÍCH VÀ DỰ BÁO TỶ GIÁ USD/VND
  \end{beamercolorbox}%
  \begin{beamercolorbox}[wd=.13\paperwidth,ht=2.65ex,dp=1.05ex,center]{date in head/foot}%
    \color{white}\scriptsize\bfseries \insertframenumber/\inserttotalframenumber
  \end{beamercolorbox}%
  }%
  }%
}
\setbeamertemplate{itemize/enumerate body begin}{\small}
\setbeamertemplate{itemize/enumerate subbody begin}{\footnotesize}
\newcommand{\topic}[1]{\vspace{0.04cm}{\color{mainred}\bfseries\small #1}\par\vspace{0.08cm}}
\newcommand{\src}[1]{\gdef\slideRef{#1}}
\newcommand{\takeaway}[1]{\vspace{0.06cm}\begin{beamercolorbox}[wd=\linewidth,sep=0.13cm,rounded=true]{block body}{\small\color{darkred}\textbf{Kết luận:} #1}\end{beamercolorbox}}
\newcommand{\notegray}[1]{{\small\color{midgray}#1}}
\newcommand{\kpi}[3]{\begin{minipage}[c]{0.30\linewidth}\centering\begin{beamercolorbox}[wd=\linewidth,sep=0.12cm,center,rounded=true]{block body}{\Large\bfseries\color{mainred}#1}\par\vspace{0.03cm}{\footnotesize\color{darkgray}#2}\par{\scriptsize\color{midgray}#3}\end{beamercolorbox}\end{minipage}}
\newcommand{\twoplot}[4]{%
\begin{columns}[T,onlytextwidth]
\begin{column}{0.49\textwidth}\centering\includegraphics[width=\linewidth,height=3.0cm,keepaspectratio]{assets/#1.png}\\[-2pt]{\scriptsize\color{midgray}#2}\end{column}
\begin{column}{0.49\textwidth}\centering\includegraphics[width=\linewidth,height=3.0cm,keepaspectratio]{assets/#3.png}\\[-2pt]{\scriptsize\color{midgray}#4}\end{column}
\end{columns}}
\newcommand{\bigfig}[3]{\centering\includegraphics[width=#2,height=#3,keepaspectratio]{assets/#1.png}}
\tikzset{flow/.style={draw=mainred!60,fill=softred,rounded corners=3pt,align=center,minimum height=0.72cm,inner sep=4pt},
arrow/.style={-{Latex[length=2.4mm]},thick,draw=mainred},dim/.style={draw=midgray!45,fill=lightgray,rounded corners=3pt,align=center,inner sep=6pt}}
\tikzset{algstep/.style={draw=mainred!65,fill=softred,rounded corners=3pt,
  text width=2.25cm,minimum height=1.02cm,align=center,inner sep=3pt,font=\small}}
\newcommand{\procflow}[8]{%
\begin{tikzpicture}[node distance=0.26cm]
\node[algstep](stepa){{\color{mainred}\bfseries\scriptsize BƯỚC 01}\\[-2pt]\textbf{#1}\\[-2pt]{\scriptsize #2}};
\node[algstep,right=of stepa](stepb){{\color{mainred}\bfseries\scriptsize BƯỚC 02}\\[-2pt]\textbf{#3}\\[-2pt]{\scriptsize #4}};
\node[algstep,right=of stepb](stepc){{\color{mainred}\bfseries\scriptsize BƯỚC 03}\\[-2pt]\textbf{#5}\\[-2pt]{\scriptsize #6}};
\node[algstep,right=of stepc](stepd){{\color{mainred}\bfseries\scriptsize BƯỚC 04}\\[-2pt]\textbf{#7}\\[-2pt]{\scriptsize #8}};
\draw[arrow](stepa)--(stepb); \draw[arrow](stepb)--(stepc); \draw[arrow](stepc)--(stepd);
\end{tikzpicture}\par%
}
\hypersetup{pdftitle={Phan tich va du bao ty gia USD/VND - Chuyen de 2, 47 slides, seven section dividers},pdfauthor={Nhom Chuyen de 2},colorlinks=false}


% ---------- Editable academic seven-part dividers and contents ----------
% All new slides are vector-based with local 16x9cm coordinates.
% They are not flattened screenshot images and can be edited in this file.
\newcommand{\specialpagebase}{%
  \path[use as bounding box] (0,0) rectangle (16,9);
  \fill[paperbg] (0,0) rectangle (16,9);
  \begin{scope}
   \clip (0,.48) rectangle (16,8.04);
   \fill[softred] (14.50,-.1) circle (2.87);
  \end{scope}
  \path[shade,left color=darkred,right color=mainred] (0,8.04) rectangle (16,9);
  % Segmented Beamer footer, matching existing pages and user-provided theme.
  \fill[darkred] (0,0) rectangle (3.68,.43);
  \fill[footerred] (3.68,0) rectangle (13.92,.43);
  \fill[mainred] (13.92,0) rectangle (16,.43);
  \node[anchor=west,text=white,inner sep=0pt] at (.31,.25)
    {{\fontsize{9}{11}\selectfont\bfseries CHUYÊN ĐỀ 2}};
  \node[anchor=center,text=white,inner sep=0pt] at (8,.25)
    {{\fontsize{9}{11}\selectfont PHÂN TÍCH VÀ DỰ BÁO TỶ GIÁ USD/VND}};
  \node[anchor=center,text=white,inner sep=0pt] at (15.05,.25)
    {{\fontsize{8}{10}\selectfont\bfseries\insertframenumber/\inserttotalframenumber}};
}
\newcommand{\partdivider}[4]{%
  \begin{frame}[plain,label=part#1]
    \noindent\hspace*{-0.72cm}%
    \begin{tikzpicture}[x=1cm,y=1cm]
      \specialpagebase
      \node[anchor=west,text=white,inner sep=0pt] at (.66,8.40)
        {{\fontsize{24}{29}\selectfont\bfseries PHẦN #1}};
      \draw[mainred,line width=1.45pt] (4.78,2.78)--(4.78,6.2);
      \node[anchor=center,text=darkred,inner sep=0pt,text width=4.38cm,align=center]
        at (2.47,4.52) {{\fontsize{75}{84}\selectfont\bfseries #1}};
      \node[anchor=west,align=left,inner sep=0pt,text width=10.18cm]
        at (5.42,4.53)
        {{\fontsize{22}{26}\selectfont\color{darkred}\bfseries #2}\par\vspace{0.16cm}
         {\fontsize{11.8}{16}\selectfont\color{darkgray}#3}};
    \end{tikzpicture}%
  \end{frame}%
}
\newcommand{\tocentry}[4]{%
  \node[anchor=north west,text=darkred,inner sep=0pt] at (#1,#2)
    {{\fontsize{26}{29}\selectfont\bfseries #3.}};
  \node[anchor=north west,inner sep=0pt,text=darkgray,align=left,text width=5.76cm]
    at ({#1+1.64},#2)
    {{\fontsize{11.9}{16.1}\selectfont\bfseries #4}};
}
\newcommand{\presentationContents}{%
  \begin{frame}[plain,label=contents]
  \noindent\hspace*{-0.72cm}%
  \begin{tikzpicture}[x=1cm,y=1cm]
     \specialpagebase
     \node[anchor=west,text=white,inner sep=0pt] at (.65,8.39)
       {{\fontsize{24}{27}\selectfont\bfseries NỘI DUNG THUYẾT TRÌNH}};
     \tocentry{.73}{6.91}{I}{Tổng quan và dữ liệu}
     \tocentry{.73}{5.51}{II}{Mô hình ARIMA}
     \tocentry{.73}{4.11}{III}{Mô hình ARMA--GARCH}
     \tocentry{.73}{2.71}{IV}{Không gian trạng thái và\\bộ lọc Kalman}
     \tocentry{8.48}{6.91}{V}{Mô hình SETAR/TAR}
     \tocentry{8.48}{5.51}{VI}{{\fontsize{11.2}{15.2}\selectfont Mô hình Markov Switching}}
     \tocentry{8.48}{4.11}{VII}{So sánh và kết luận}
  \end{tikzpicture}%
  \end{frame}%
}
\newcommand{\presentationThankYou}{%
  \begin{frame}[plain,label=thanks]
  \noindent\hspace*{-0.72cm}%
  \begin{tikzpicture}[x=1cm,y=1cm]
     \specialpagebase
     \node[anchor=center,text=white,inner sep=0pt] at (8,8.40)
       {{\fontsize{24}{28}\selectfont\bfseries CHUYÊN ĐỀ 2}};
     \node[anchor=center,text=darkred,align=center,inner sep=0pt,text width=14.6cm]
       at (8,5.22)
       {{\fontsize{22.2}{27}\selectfont\bfseries XIN TRÂN TRỌNG CẢM ƠN!}};
     \node[anchor=center,text=darkgray,align=center,inner sep=0pt,text width=13.6cm]
       at (8,4.08)
       {{\fontsize{13.7}{20}\selectfont Nhóm mong nhận được câu hỏi và góp ý từ giảng viên.}};
     \node[anchor=center,text=darkred,inner sep=0pt] at (8,2.50)
       {{\fontsize{14}{18}\selectfont Phân tích và dự báo tỷ giá USD/VND}};
  \end{tikzpicture}%
  \end{frame}%
}

\begin{document}

% SLIDE 01 -----------------------------------------------------------------
\begin{frame}[plain]
\begin{tikzpicture}[remember picture,overlay]
  \fill[mainred] (current page.north west) rectangle ([xshift=0.23cm]current page.south west);
  \fill[mainred] ([yshift=-2.38cm]current page.north west) rectangle ([yshift=-2.42cm]current page.north east);
  \fill[softred] ([xshift=-3.2cm,yshift=0.5cm]current page.south east) circle (2.3cm);
\end{tikzpicture}
\vspace{0.10cm}\centering
{\footnotesize\bfseries ĐẠI HỌC QUỐC GIA HÀ NỘI}\\[1pt]
{\footnotesize\bfseries TRƯỜNG ĐẠI HỌC KHOA HỌC TỰ NHIÊN}\\[1pt]
{\scriptsize KHOA TOÁN -- CƠ -- TIN HỌC}\\[0.60cm]
{\color{mainred}\bfseries\large BÁO CÁO TIỂU LUẬN -- CHUYÊN ĐỀ 2}\\[0.21cm]
{\color{darkred}\bfseries\LARGE PHÂN TÍCH VÀ DỰ BÁO}\\[0.08cm]
{\color{darkred}\bfseries\LARGE TỶ GIÁ USD/VND}\\[0.25cm]
{\large Bằng năm mô hình chuỗi thời gian}\\[0.34cm]
{\small Giảng viên: \textbf{TS. Hoàng Thị Phương Thảo}}\\[0.12cm]
{\footnotesize Trần Tùng Lâm \enspace $\bullet$ \enspace Nguyễn Thị Nhật Linh \enspace $\bullet$ \enspace Nguyễn Phúc Lợi}\\[0.05cm]
{\footnotesize Hoàng Thanh Tùng \enspace $\bullet$ \enspace Tống Việt Tú}\\[0.21cm]
{\small\color{midgray}Hà Nội, 2026}
\end{frame}

% Redesigned TOC based on the supplied reference.
\presentationContents

\section{I. Tổng quan và dữ liệu}
\partdivider{I}{TỔNG QUAN VÀ DỮ LIỆU}{{\fontsize{10.5}{14.7}\selectfont Bài toán nghiên cứu, dữ liệu USD/VND và thiết kế đánh giá}}{}

% SLIDE 02 -----------------------------------------------------------------
\begin{frame}{Bài toán và năm cách tiếp cận}
\topic{Câu hỏi nghiên cứu}
\centering\large\(\boxed{\mathcal F_t\ \longrightarrow\ \widehat P_{t+1\mid t}}\)\\[0.14cm]
\notegray{Dùng lịch sử giá đóng cửa tới hết phiên $t$ để dự báo giá đóng cửa phiên kế tiếp.}\\[0.25cm]
\begin{tikzpicture}[node distance=0.13cm]
\node[flow,text width=2.15cm] (a) {\textbf{ARIMA}\\\scriptsize Phụ thuộc tuyến tính};
\node[flow,text width=2.15cm,right=of a] (b) {\textbf{ARMA--GARCH}\\\scriptsize Phương sai thay đổi};
\node[flow,text width=2.15cm,right=of b] (c) {\textbf{Kalman}\\\scriptsize Trạng thái liên tục ẩn};
\node[flow,text width=2.15cm,right=of c] (d) {\textbf{SETAR}\\\scriptsize Ngưỡng quan sát};
\node[flow,text width=2.15cm,right=of d] (e) {\textbf{Markov}\\\scriptsize Trạng thái rời rạc ẩn};
\end{tikzpicture}
\vspace{0.38cm}
\begin{beamercolorbox}[wd=\linewidth,sep=0.19cm,rounded=true]{block body}
\small \textbf{Đối chứng:} Random Walk dự báo $\widehat P_{t+1\mid t}=P_t$; \textbf{không phải mô hình thứ sáu}. Mục tiêu là kiểm tra liệu năm đặc tả phức tạp có cải thiện \textbf{dự báo ngoài mẫu} hay không.
\end{beamercolorbox}
\src{Phần I, mục 1, 4.6; Phần II, mục 1}
\end{frame}

% Phần I tiếp tục: dữ liệu và phương pháp

% SLIDE 03 -----------------------------------------------------------------
\begin{frame}{Dữ liệu USD/VND và chất lượng nguồn}
\vspace{0.16cm}\centering
\kpi{2.784}{Phiên quan sát}{04/01/2016--31/08/2026}\hfill
\kpi{569}{Dòng bất nhất OHLC}{Đánh dấu, chưa tự sửa}\hfill
\kpi{1.777}{Giá trị Volume thiếu}{Không dùng làm đầu vào}\par\vspace{0.34cm}
\begin{columns}[T,onlytextwidth]
\begin{column}{0.52\linewidth}
\topic{Tệp đầu vào}
\begin{tabular}{@{}ll@{}}\toprule
\textbf{Nhóm trường} & \textbf{Tên chuẩn hóa}\\\midrule
Ngày & Date\\
Giá & Open, High, Low, Close\\
Giao dịch & Volume\\
Thay đổi & Change\_pct\\\bottomrule
\end{tabular}
\end{column}
\begin{column}{0.45\linewidth}
\small
\begin{itemize}\setlength{\itemsep}{6pt}
\item Biến mục tiêu là \textbf{Close} (VND/USD).
\item \textbf{Không thiếu Close}, không trùng ngày.
\item Nguồn dữ liệu: Investing.com.
\end{itemize}
\end{column}
\end{columns}
\src{Phần II, mục 2.1--2.3, Bảng TN.2--TN.3}
\end{frame}

% SLIDE 04 -----------------------------------------------------------------
\begin{frame}{Tiền xử lý và hai chuỗi làm việc}
\vspace{0.09cm}
\begin{tikzpicture}[every node/.style={font=\small}]
\node[flow,text width=1.80cm] (raw) at (0.95,0) {\textbf{CSV gốc}\\2.784 phiên};
\node[flow,text width=2.20cm] (price) at (4.15,0) {\textbf{Chuẩn hóa}\\Ngày, giá, sắp phiên};
\node[flow,text width=2.55cm] (lp) at (8.40,.68) {\textbf{log\_price}\\ARIMA, Kalman};
\node[flow,text width=2.55cm] (ret) at (8.40,-.72) {\textbf{log\_return}\\GARCH, SETAR, Markov};
\draw[arrow] (raw)--(price);
\draw[arrow] (price.east)--++(.28,0)|-(lp.west);
\draw[arrow] (price.east)--++(.28,0)|-(ret.west);
\end{tikzpicture}
\vspace{0.3cm}
\begin{columns}[T,onlytextwidth]
\begin{column}{.48\linewidth}
\[y_t=\ln P_t\qquad r_t=100(\ln P_t-\ln P_{t-1})\]
\end{column}
\begin{column}{.48\linewidth}
\[\widehat P_{t+1}=P_t\exp\!\left(\frac{\widehat r_{t+1}}{100}\right)\]
\end{column}
\end{columns}
\small\begin{itemize}\setlength{\itemsep}{4pt}
\item Giữ nguyên lịch phiên; không chèn ngày nghỉ; loại phiên đầu \emph{chỉ trong chuỗi lợi suất}.
\item Tạo 2.784 quan sát giá và 2.783 quan sát lợi suất. Mọi dự báo được so trên \textbf{VND/USD}.
\end{itemize}
\src{Phần I, mục 1.2; Phần II, mục 2.2--2.6}
\end{frame}

% SLIDE 05 -----------------------------------------------------------------
\begin{frame}{EDA: xu thế tỷ giá và biến động lợi suất}
\twoplot{fig01_close}{Hình 1 -- Tỷ giá Close, train 2016--2024}{fig02_returns}{Hình 2 -- Lợi suất logarit, train 2016--2024}
\vspace{0.19cm}
\begin{columns}[T,onlytextwidth]
\begin{column}{0.66\linewidth}
\small\begin{itemize}\setlength{\itemsep}{3pt}
\item Close có xu thế và tính bền vững rõ rệt.
\item Lợi suất dao động quanh 0, có cụm biến động lớn và quan sát cực trị.
\end{itemize}
\end{column}
\begin{column}{0.32\linewidth}
\begin{beamercolorbox}[wd=\linewidth,sep=.15cm,rounded=true]{block body}
\centering{\Large\color{mainred}\bfseries 17,475}\\[-3pt]{\scriptsize Excess kurtosis lợi suất train}
\end{beamercolorbox}
\end{column}
\end{columns}
\src{Phần II, mục 4, Bảng TN.7, Hình 1--3}
\end{frame}

% SLIDE 06 -----------------------------------------------------------------
\begin{frame}{Các kiểm định tiền mô hình hóa}
\vspace{0.12cm}
{\small\renewcommand{\arraystretch}{1.35}
\begin{tabularx}{\linewidth}{@{}>{\bfseries}p{2.8cm}p{2.8cm}X@{}}
\rowcolor{softred}\textbf{Kiểm định} & \textbf{Kết quả train} & \textbf{Hàm ý trong nghiên cứu}\\\midrule
ADF: log\_price & $p\approx0{,}908$ & Chưa bác bỏ nghiệm đơn vị; xem xét sai phân $d=1$.\\
ADF: log\_return & $p\approx1{,}815\cdot10^{-28}$ & Bác bỏ nghiệm đơn vị theo đặc tả ADF đã dùng.\\
Ljung--Box & $p$ rất nhỏ & Lợi suất có phụ thuộc ở các độ trễ xét.\\
ARCH--LM(10) & $p\approx1{,}049\cdot10^{-53}$ & Xem xét phương sai có điều kiện thay đổi.\\\bottomrule
\end{tabularx}}
\vspace{0.26cm}
\takeaway{ADF và ARCH--LM định hướng thiết kế; \emph{không tự chứng minh} tính dừng tuyệt đối hay hiệu quả dự báo của một mô hình.}
\src{Phần II, mục 5, Bảng TN.9--TN.11}
\end{frame}

% SLIDE 07 -----------------------------------------------------------------
\begin{frame}{Chia dữ liệu theo thời gian và dự báo cuốn chiếu}
\vspace{0.14cm}
\begin{tikzpicture}[x=1cm,y=1cm]
\fill[mainred] (0,0) rectangle (7.0,1.02);
\fill[footerred] (7.06,0) rectangle (9.18,1.02);
\fill[darkred] (9.24,0) rectangle (11.86,1.02);
\node[white,align=center] at (3.5,0.55) {\bfseries Train\\\scriptsize 2016--2024 · 2.347 phiên};
\node[white,align=center] at (8.12,0.55) {\bfseries Validation\\\scriptsize 2025 · 262 phiên};
\node[white,align=center] at (10.56,0.55) {\bfseries Test\\[-1pt]\tiny 01--08/2026\\[-2pt]\tiny 175 phiên};
\end{tikzpicture}\par
\vspace{0.34cm}
\topic{Tại mỗi phiên test: dự báo phải đi trước dữ liệu thực tế}
\centering
\procflow
{Thông tin}{Có đến phiên $t$}
{Dự báo}{Giá phiên $t+1$}
{Quan sát}{Giá thực $P_{t+1}$}
{Cập nhật}{Bộ lọc / phần dư}
\vspace{0.26cm}
\notegray{Dự báo một bước với gốc dịch chuyển, không phải dự báo cả 175 phiên từ một mốc cố định. Cấu hình không được chọn lại theo test.}
\src{Phần II, mục 3, Bảng TN.6}
\end{frame}

% SLIDE 08 -----------------------------------------------------------------
\begin{frame}{Random Walk và tiêu chí đánh giá}
\topic{Mốc đối chứng bắt buộc}
\centering\[\boxed{\widehat P^{\text{RW}}_{t+1\mid t}=P_t}\]
\vspace{0.10cm}
\kpi{18,56286}{MAE test}{VND/USD}\hspace{0.35cm}\kpi{30,05340}{RMSE test}{VND/USD}\par\vspace{0.22cm}
\begin{columns}[T,onlytextwidth]
\begin{column}{0.51\linewidth}
\[\mathrm{MAE}=\frac1n\sum_{t=1}^n|e_t|\]
\end{column}
\begin{column}{0.45\linewidth}
\[\mathrm{RMSE}=\sqrt{\frac1n\sum_{t=1}^n e_t^2}\]
\end{column}
\end{columns}
\vspace{0.10cm}
\small\centering $\mathrm{RMSE\ ratio}=\mathrm{RMSE}_{\mathrm{model}}/\mathrm{RMSE}_{\mathrm{RW}}$\par
\takeaway{RMSE thấp hơn về số học không tự chứng minh khác biệt có ý nghĩa thống kê.}
\vspace{-0.13cm}
\src{Phần II, mục 6 và 8.2, Bảng TN.12}
\end{frame}

\section{II. Mô hình ARIMA}
\partdivider{II}{MÔ HÌNH ARIMA}{Cơ sở lý thuyết và kết quả thực nghiệm}{}

% SLIDE 09 -----------------------------------------------------------------
\begin{frame}{ARIMA: khái niệm và phương trình}
\vspace{0.17cm}
\centering\[\boxed{\phi(B)(1-B)^d y_t=c+\theta(B)\varepsilon_t}\]
\vspace{0.12cm}
\begin{columns}[T,onlytextwidth]
\begin{column}{.32\linewidth}
\begin{block}{AR($p$)}\small Phụ thuộc \textbf{tuyến tính} vào $p$ giá trị trễ.\end{block}
\end{column}
\begin{column}{.32\linewidth}
\begin{block}{I($d$)}\small Sai phân cấp $d$ để xét chuỗi có tính dừng.\end{block}
\end{column}
\begin{column}{.32\linewidth}
\begin{block}{MA($q$)}\small Phụ thuộc vào $q$ \textbf{nhiễu đổi mới} trễ.\end{block}
\end{column}
\end{columns}
\vspace{0.24cm}
\small $B$ là toán tử trễ, $y_t=\ln P_t$; kiểm định ADF trên dữ liệu USD/VND gợi ý xét $d=1$.
\takeaway{Điều kiện nhân quả/khả nghịch và chẩn đoán phần dư là điều kiện kỹ thuật; không bảo đảm dự báo ngoài mẫu tốt.}
\src{Phần I, mục 4.1.1--4.1.5}
\end{frame}

% SLIDE 10 -----------------------------------------------------------------
\begin{frame}{ARIMA: Box--Jenkins và dự báo một bước}
\vspace{0.18cm}
\centering
\procflow
{Nhận dạng}{ADF, ACF/PACF}
{Ước lượng}{MLE; AIC/BIC}
{Chẩn đoán}{Hội tụ; Ljung--Box}
{Lựa chọn}{RMSE trên val}
\vspace{0.41cm}
\topic{Dự báo trên thang logarit khi $d=1$}
\centering\[\widehat y_{t+1\mid t}=y_t+\widehat{\Delta y}_{t+1\mid t}\]
\small\begin{itemize}\setlength{\itemsep}{5pt}
\item Chọn $p,q$ từ ứng viên hội tụ, thỏa điều kiện mô hình và có sai số val phù hợp.
\item Chuyển dự báo về giá gốc trước khi tính MAE/RMSE.
\end{itemize}
\takeaway{AIC/BIC không thay thế kiểm tra phần dư và sai số ngoài mẫu.}
\vspace{-0.13cm}
\src{Phần I, mục 4.1.6--4.1.8; Phần II, mục 7.1}
\end{frame}

% SLIDE 11 -----------------------------------------------------------------
\begin{frame}{ARIMA: lựa chọn cấu hình (3,1,1)}
\begin{columns}[T,onlytextwidth]
\begin{column}{.62\linewidth}
\bigfig{p1_2_validation}{\linewidth}{4.17cm}
\centering{\scriptsize\color{midgray}Hình P1.2 -- RMSE validation của ứng viên hợp lệ}
\end{column}
\begin{column}{.35\linewidth}
\topic{Quy tắc lựa chọn}
\small\begin{enumerate}\setlength{\itemsep}{7pt}
\item Thử $25$ cấu hình $p,q\in\{0,\dots,4\}$.
\item Loại mô hình không hợp lệ / không hội tụ.
\item Chấp nhận RMSE val trong vùng \textbf{1\%} của tốt nhất, ưu tiên ít tham số.
\end{enumerate}
\vspace{0.1cm}
\begin{beamercolorbox}[wd=\linewidth,sep=.13cm,rounded=true]{block body}
\centering\textbf{ARIMA(3,1,1)}\\[3pt]{\color{mainred}\bfseries RMSE val $\approx43{,}1253$}
\end{beamercolorbox}
\end{column}
\end{columns}
\src{Phần II, mục 7.1.2--7.1.4, Hình P1.2, Bảng P1.1--P1.3}
\end{frame}

% SLIDE 12 -----------------------------------------------------------------
\begin{frame}{ARIMA: dự báo test và khoảng 95\%}
\begin{columns}[T,onlytextwidth]
\begin{column}{.68\linewidth}
\bigfig{p1_3_forecast}{\linewidth}{4.25cm}
\centering{\scriptsize\color{midgray}Hình P1.3 -- Thực tế, dự báo một bước và khoảng 95\%}
\end{column}
\begin{column}{.29\linewidth}
\topic{Tập test: 175 phiên}
\small\begin{tabular}{@{}lr@{}}\toprule
& \textbf{ARIMA}\\\midrule
MAE & 18,7363\\
RMSE & 30,4392\\
Bao phủ 95\% & 95,43\%\\\bottomrule
\end{tabular}
\vspace{0.27cm}
\begin{beamercolorbox}[wd=\linewidth,sep=.13cm,rounded=true]{block body}
\small RMSE \textbf{cao hơn RW khoảng 1,28\%} (RW: 30,0534).
\end{beamercolorbox}
\end{column}
\end{columns}
\takeaway{Khoảng dự báo có độ bao phủ gần danh nghĩa, nhưng dự báo điểm không vượt Random Walk.}
\src{Phần II, mục 7.1.5--7.1.6, Bảng P1.4--P1.5}
\end{frame}

% SLIDE 13 -----------------------------------------------------------------
\begin{frame}{ARIMA: chẩn đoán và độ vững}
\begin{columns}[T,onlytextwidth]
\begin{column}{.57\linewidth}
\bigfig{p1_4_diagnostic}{\linewidth}{3.87cm}
\centering{\scriptsize\color{midgray}Hình P1.4 -- Phần dư và phân phối sai số}
\end{column}
\begin{column}{.40\linewidth}
\topic{Kiểm tra sau ước lượng}
\small\begin{itemize}\setlength{\itemsep}{8pt}
\item Phần dư train còn đuôi dày / hiệu ứng ARCH.
\item Sai số dự báo test còn tự tương quan.
\item HAC trên chênh lệch loss bình phương với RW: $p\approx0{,}469$.
\item Bootstrap 95\% của $\Delta\mathrm{RMSE}$ chứa 0.
\end{itemize}
\end{column}
\end{columns}
\takeaway{ARIMA nắm bắt quan hệ tuyến tính ngắn hạn nhưng chưa chứng minh được lợi thế dự báo một bước so với RW.}
\src{Phần II, mục 7.1.6--7.1.10, Bảng P1.6--P1.8}
\end{frame}

\section{III. Mô hình ARMA--GARCH}
\partdivider{III}{MÔ HÌNH ARMA--GARCH}{Trung bình, phương sai có điều kiện và dự báo}{}

% SLIDE 14 -----------------------------------------------------------------
\begin{frame}{ARMA--GARCH: trung bình và phương sai}
\topic{Hai thành phần được mô hình hóa đồng thời}
\centering
\begin{columns}[T,onlytextwidth]
\begin{column}{.48\linewidth}
\begin{block}{Trung bình có điều kiện (ARMA)}
\[r_t=\mu_t+\varepsilon_t\]
\centering\small Dự báo \textbf{lợi suất}.
\end{block}
\end{column}
\begin{column}{.48\linewidth}
\begin{block}{Phương sai có điều kiện (GARCH)}
\[\varepsilon_t=\sqrt{h_t}\,z_t\]
\[h_t=\omega+\alpha\varepsilon_{t-1}^2+\beta h_{t-1}\]
\centering\small Dự báo \textbf{độ biến động}.
\end{block}
\end{column}
\end{columns}
\vspace{0.28cm}
\small\begin{itemize}\setlength{\itemsep}{6pt}
\item $\omega$: mức nền; $\alpha$: tác động cú sốc mới; $\beta$: mức duy trì của biến động.
\item Với GARCH(1,1), $\omega>0$, $\alpha,\beta\geq0$ và $\alpha+\beta<1$ là điều kiện quen thuộc cho phương sai vô điều kiện hữu hạn dưới giả định phù hợp.
\end{itemize}
\takeaway{Mô tả rủi ro/biến động tốt hơn không tự động tạo dự báo giá chính xác hơn.}
\vspace{-0.13cm}
\src{Phần I, mục 4.2.1--4.2.3}
\end{frame}

% SLIDE 15 -----------------------------------------------------------------
\begin{frame}{ARMA--GARCH: ước lượng và dự báo phân phối}
\vspace{0.06cm}
\procflow
{Ước lượng}{ARMA + GARCH}
{Phân phối}{Student--$t$}
{Phân vị}{Khoảng 95\%}
{Chuyển đổi}{VND/USD}
\vspace{0.31cm}
\centering\[\boxed{\widehat P^{\mathrm{median}}_{t+1}=P_t\exp\!\left(\frac{\widehat r_{t+1}}{100}\right)}\]
\vspace{0.10cm}
\small\begin{itemize}\setlength{\itemsep}{6pt}
\item Student--$t$ chuẩn hóa cho phép nhiễu đuôi dày; \emph{không} tự áp dụng công thức kỳ vọng log-chuẩn cho Student--$t$.
\item Đánh giá điểm bằng MAE/RMSE; đánh giá biến động bằng QLIKE trên đại lượng thay thế và kiểm tra khoảng/VaR.
\end{itemize}
\src{Phần I, mục 4.2.4--4.2.6}
\end{frame}

% SLIDE 16 -----------------------------------------------------------------
\begin{frame}{ARMA--GARCH: cấu hình và tham số cuối}
\begin{columns}[T,onlytextwidth]
\begin{column}{.54\linewidth}
\topic{Chọn trên validation}
\small\begin{itemize}\setlength{\itemsep}{5pt}
\item 4 cấu trúc ARMA $\times$ 2 phân phối nhiễu = \textbf{8 ứng viên}.
\item Cố định GARCH(1,1); chọn theo \textbf{MAE của giá} trên val.
\item Đặc tả được chọn: \textbf{ARMA(0,0)--GARCH(1,1), Student--$t$}.
\end{itemize}
\end{column}
\begin{column}{.43\linewidth}
{\small\renewcommand{\arraystretch}{1.27}\begin{tabular}{@{}lr@{}}\toprule
Tham số & Ước lượng cuối\\\midrule
$\mu$ & $-0{,}000118$\\
$\omega$ & $0{,}000072$\\
$\alpha$ & $0{,}201396$\\
$\beta$ & $0{,}796878$\\
$\nu$ & $3{,}404532$\\
$\alpha+\beta$ & $0{,}998273$\\\bottomrule
\end{tabular}}
\end{column}
\end{columns}
\vspace{0.08cm}
\centering\[r_t\approx-0{,}000118+\varepsilon_t,\quad h_t=0{,}000072+0{,}201396\varepsilon_{t-1}^2+0{,}796878h_{t-1}\]
\takeaway{$\alpha+\beta$ rất sát 1 và gần ràng buộc tối ưu; thận trọng khi diễn giải tính dai dẳng của phương sai.}
\vspace{-0.13cm}
\src{Phần II, mục 7.2.2--7.2.3, Bảng P2.1--P2.2}
\end{frame}

% SLIDE 17 -----------------------------------------------------------------
\begin{frame}{ARMA--GARCH: dự báo giá và độ biến động}
\begin{columns}[T,onlytextwidth]
\begin{column}{.65\linewidth}
\bigfig{p2_2_garch_pi}{\linewidth}{4.14cm}
\centering{\scriptsize\color{midgray}Hình P2.2 -- Khoảng giá 95\% và độ lệch chuẩn GARCH}
\end{column}
\begin{column}{.32\linewidth}
\topic{Test: 175 phiên}
\small\begin{tabular}{@{}lr@{}}\toprule
MAE & 18,56303\\
RMSE & 30,05206\\
RW RMSE & 30,05340\\
Bao phủ & 93,714\%\\
Độ rộng & 94,53\\\bottomrule
\end{tabular}
\vspace{0.23cm}
\small \textbf{QLIKE} (trên đại lượng thay thế):\\
GARCH: $-3{,}80869$\\
Phương sai cố định: $-3{,}26101$.
\end{column}
\end{columns}
\takeaway{Dự báo điểm gần như bằng RW; mô hình vẫn cung cấp dự báo phương sai biến thiên theo thời gian.}
\src{Phần II, mục 7.2.5--7.2.6, Bảng P2.4--P2.5}
\end{frame}

% SLIDE 18 -----------------------------------------------------------------
\begin{frame}{ARMA--GARCH: VaR và chẩn đoán}
\topic{Kiểm tra VaR đuôi trái ở mức 5\%}
\centering
\kpi{5\%}{Tần suất danh nghĩa}{Kỳ vọng theo mô hình}\hfill
\kpi{15/175}{Số phiên vi phạm}{Tương ứng 8,571\%}\hfill
\kpi{0,03685}{$p$ nhị thức hai phía}{Dưới giả định độc lập}\par\vspace{0.34cm}
\begin{columns}[T,onlytextwidth]
\begin{column}{.48\linewidth}
\begin{block}{Phương trình trung bình}\small Ljung--Box trên phần dư chuẩn hóa còn bác bỏ giả thuyết không tự tương quan ở một số độ trễ.\end{block}
\end{column}
\begin{column}{.48\linewidth}
\begin{block}{Phương trình phương sai}\small ARCH--LM và Ljung--Box bình phương phần dư chưa ghi nhận ARCH còn sót tại độ trễ đã xét.\end{block}
\end{column}
\end{columns}
\takeaway{Khoảng hẹp và QLIKE thấp hơn không đồng nghĩa ngưỡng rủi ro đuôi được hiệu chỉnh tốt; chưa có kiểm định loss ghép cặp về dự báo giá so với RW.}
\src{Phần II, mục 7.2.4, 7.2.7--7.2.8, Bảng P2.3, P2.6}
\end{frame}

\section{IV. Không gian trạng thái và Kalman}
\partdivider{IV}{{\fontsize{19}{24}\selectfont KHÔNG GIAN TRẠNG THÁI\\VÀ BỘ LỌC KALMAN}}{Trạng thái ẩn, lọc và dự báo một bước}{}

% SLIDE 19 -----------------------------------------------------------------
\begin{frame}{Mô hình không gian trạng thái: cấu trúc ẩn}
\topic{Hai phương trình của mô hình tuyến tính Gaussian}
\centering
\[\boxed{\begin{aligned}y_t&=Z_t\alpha_t+\varepsilon_t, &\varepsilon_t&\sim N(0,H_t)\\
\alpha_{t+1}&=T_t\alpha_t+\eta_t, &\eta_t&\sim N(0,Q_t)
\end{aligned}}\]
\vspace{0.12cm}
\begin{columns}[T,onlytextwidth]
\begin{column}{.52\linewidth}
\topic{Ý nghĩa}
\small\begin{itemize}\setlength{\itemsep}{5pt}
\item $\alpha_t$ là trạng thái ẩn: mức hoặc mức + xu hướng.
\item $H_t$: nhiễu quan sát; $Q_t$: nhiễu của quá trình chuyển trạng thái.
\end{itemize}
\end{column}
\begin{column}{.45\linewidth}
\topic{Ba cấu trúc khảo sát}
\small\begin{tabular}{@{}ll@{}}\toprule
Local Level & Mức ẩn\\
Local Linear Trend & Mức + độ dốc\\
Smooth Trend & Xu hướng trơn\\\bottomrule
\end{tabular}
\end{column}
\end{columns}
\takeaway{Không thể kết luận mô hình nhiều trạng thái ẩn hơn sẽ cho dự báo tốt hơn.}
\src{Phần I, mục 4.3.1--4.3.2, Bảng LT.1}
\end{frame}

% SLIDE 20 -----------------------------------------------------------------
\begin{frame}{Bộ lọc Kalman: bốn bước cập nhật}
\vspace{0.08cm}
\centering
\procflow
{Dự đoán}{$a^-_t$, $C^-_t$}
{Đổi mới}{$v_t=y_t-Z_ta^-_t$}
{Tính Gain}{Ma trận $K_t$}
{Cập nhật}{$a_t^+$, $C_t^+$}
\vspace{0.42cm}
\[\boxed{K_t=C_t^-Z_t^{\mathsf T}\!\left(Z_tC_t^-Z_t^{\mathsf T}+H_t\right)^{-1}}\]
\[\boxed{a_t^+=a_t^-+K_t\bigl(y_t-Z_ta_t^-\bigr)}\]
\vspace{0.11cm}
\small $K_t$ phân bổ trọng số giữa ước lượng trước cập nhật và thông tin mới. Trong Local Level vô hướng, $K_t\to1$ khi nhiễu quan sát $H$ rất nhỏ so với độ bất định trạng thái.
\src{Phần I, mục 4.3.3--4.3.4}
\end{frame}

% SLIDE 21 -----------------------------------------------------------------
\begin{frame}{Kalman: lọc, làm trơn và dự báo ngoài mẫu}
\vspace{0.18cm}
\begin{columns}[T,onlytextwidth]
\begin{column}{.48\linewidth}
\begin{block}{Filtered / predicted state}
\small $a_{t\mid t}$ chỉ dùng $y_1,\ldots,y_t$; từ đây lập dự báo cho phiên $t+1$. \textbf{Được phép} sử dụng khi đánh giá ngoài mẫu.
\end{block}
\end{column}
\begin{column}{.48\linewidth}
\begin{block}{Smoothed state}
\small $a_{t\mid T}$ dùng cả $y_{t+1},\ldots,y_T$. Phù hợp nghiên cứu lịch sử, \textbf{không được dùng} để giả lập dự báo tại $t$.
\end{block}
\end{column}
\end{columns}
\vspace{0.20cm}
\topic{Chuyển dự báo logarit sang kỳ vọng giá (log-Gaussian)}
\centering\[\boxed{E(P_{t+1}\mid\mathcal F_t)=\exp\!\left(m_{t+1}+\frac12s_{t+1}^2\right)}\]
\vspace{0.15cm}
\small Trên test: tạo dự báo trước khi cập nhật trạng thái; \texttt{append(refit=False)} \textbf{không ước lượng lại tham số}.
\src{Phần I, mục 4.3.5--4.3.7; Phần II, mục 7.3.2}
\end{frame}

% SLIDE 22 -----------------------------------------------------------------
\begin{frame}{Kalman: chọn Local Level và giải thích Gain}
\begin{columns}[T,onlytextwidth]
\begin{column}{.55\linewidth}
\topic{RMSE validation theo cấu trúc}
\small\renewcommand{\arraystretch}{1.40}
\begin{tabular}{@{}lr@{}}\toprule
Cấu trúc & RMSE (VND/USD)\\\midrule
\rowcolor{softred}\textbf{Local Level} & \textbf{43,2512}\\
Local Linear Trend & 43,9451\\
Smooth Trend & 49,4510\\
Random Walk & 43,2528\\\bottomrule
\end{tabular}
\end{column}
\begin{column}{.41\linewidth}
\topic{Tham số train + val}
\small\renewcommand{\arraystretch}{1.42}
\begin{tabular}{@{}ll@{}}\toprule
$H$ & $6{,}764\cdot10^{-12}$\\
$Q$ & $1{,}893\cdot10^{-6}$\\
$H/Q$ & $3{,}573\cdot10^{-6}$\\
$K_t$ & $\approx0{,}999996$\\\bottomrule
\end{tabular}
\end{column}
\end{columns}
\vspace{0.34cm}
\begin{beamercolorbox}[wd=\linewidth,sep=.20cm,rounded=true]{block body}
\centering \textbf{$K\approx1$} $\Rightarrow$ trạng thái sau lọc gần bằng tỷ giá logarit mới quan sát $\Rightarrow$ dự báo của \textbf{Local Level gần Random Walk}.
\end{beamercolorbox}
\src{Phần II, mục 7.3.3--7.3.4, Bảng P3.1--P3.2}
\end{frame}

% SLIDE 23 -----------------------------------------------------------------
\begin{frame}{Kalman: dự báo test và kiểm định}
\begin{columns}[T,onlytextwidth]
\begin{column}{.64\linewidth}
\bigfig{p3_1_kalman_pi}{\linewidth}{4.02cm}
\centering{\scriptsize\color{midgray}Hình P3.1 -- Dự báo Local Level và khoảng 95\%}
\end{column}
\begin{column}{.33\linewidth}
\topic{Kalman / Random Walk}
\small\begin{tabular}{@{}lrr@{}}\toprule
& Kalman & RW\\\midrule
MAE &18,568&18,563\\
RMSE&30,055&30,053\\
\end{tabular}
\vspace{0.17cm}
\small Bao phủ: \textbf{94,86\%}\\
Độ rộng: \textbf{141,57} VND\\
HAC (loss bình phương): \textbf{$p\approx0{,}553$}.
\end{column}
\end{columns}
\takeaway{RMSE không có khác biệt được chứng minh theo kiểm định HAC; $\Delta$MAE dương rất nhỏ có ý nghĩa theo kiểm định riêng đã công bố.}
\src{Phần II, mục 7.3.5--7.3.10, Bảng P3.3--P3.8}
\end{frame}

\section{V. Mô hình SETAR/TAR}
\partdivider{V}{MÔ HÌNH SETAR/TAR}{Cấu trúc ngưỡng và kết quả thực nghiệm}{}

% SLIDE 24 -----------------------------------------------------------------
\begin{frame}{SETAR/TAR: hồi quy theo ngưỡng quan sát}
\topic{Hai phương trình theo dấu của lợi suất trễ so với ngưỡng}
\centering\[\boxed{r_t=\begin{cases}
a_{0,L}+\displaystyle\sum_{i=1}^p a_{i,L}r_{t-i}+\varepsilon_{t,L},& r_{t-d}\leq\gamma,\\[4pt]
a_{0,H}+\displaystyle\sum_{i=1}^p a_{i,H}r_{t-i}+\varepsilon_{t,H},& r_{t-d}>\gamma.
\end{cases}}\]
\vspace{0.15cm}
\begin{columns}[T,onlytextwidth]
\begin{column}{.48\linewidth}
\small\textbf{TAR:} chế độ do biến ngưỡng quan sát được quyết định. \textbf{SETAR:} ngưỡng là trễ của chính chuỗi đang phân tích.
\end{column}
\begin{column}{.48\linewidth}
\small $p$: bậc AR; $d$: độ trễ ngưỡng; $\gamma$: ngưỡng chia \textbf{Low / High}. Không đồng nhất Low/High với tỷ giá giảm/tăng.
\end{column}
\end{columns}
\takeaway{Khác Markov, chế độ SETAR được chọn bằng một điều kiện quan sát, không qua xác suất trạng thái ẩn.}
\vspace{-0.13cm}
\src{Phần I, mục 4.4.1--4.4.2}
\end{frame}

% SLIDE 25 -----------------------------------------------------------------
\begin{frame}{SETAR: tìm ngưỡng và thiết kế dự báo}
\vspace{0.14cm}
\centering
\procflow
{Chọn $p,d$}{Lưới trên train}
{Tìm ngưỡng}{$\gamma$: tối thiểu SSR}
{Đánh giá}{RMSE trên val}
{Ước lượng lại}{Train + val}
\vspace{0.33cm}
\small\begin{itemize}\setlength{\itemsep}{7pt}
\item Ngưỡng hợp lệ phải có đủ quan sát ở mỗi chế độ và ma trận thiết kế đủ hạng.
\item Chế độ phiên kế tiếp được xác định từ $r_{t-d}$; dự báo theo hệ số AR của chế độ đó.
\item Chuyển về mức tỷ giá và dùng \textbf{smearing theo chế độ} để giảm sai lệch do chuyển thang phi tuyến.
\end{itemize}
\takeaway{Chưa thực hiện kiểm định ngưỡng chuyên biệt; việc giảm SSR chưa chứng minh chắc chắn hiệu ứng ngưỡng tồn tại.}
\src{Phần I, mục 4.4.3--4.4.7; Phần II, mục 7.4.1}
\end{frame}

% SLIDE 26 -----------------------------------------------------------------
\begin{frame}{SETAR: cấu hình (6;4) và hai chế độ}
\begin{columns}[T,onlytextwidth]
\begin{column}{.64\linewidth}
\bigfig{p4_1_threshold}{\linewidth}{3.77cm}
\centering{\scriptsize\color{midgray}Hình P4.1 -- Hồ sơ SSR theo ngưỡng, train và train + val}
\end{column}
\begin{column}{.33\linewidth}
\topic{Đặc tả cuối}
\small\textbf{SETAR($p=6,d=4$)}\\[5pt]
Ngưỡng: $\gamma\approx -0{,}038808$ điểm \%.
\vspace{0.23cm}
\small\renewcommand{\arraystretch}{1.28}
\begin{tabular}{@{}lrr@{}}\toprule
& Low & High\\\midrule
Train + val &624&1.974\\
Test &53&122\\\bottomrule
\end{tabular}
\end{column}
\end{columns}
\takeaway{Các hệ số giữa hai chế độ khác nhau về số học, nhưng chưa có kiểm định ngưỡng riêng để xác nhận cấu trúc này.}
\src{Phần II, mục 7.4.2--7.4.3, Hình P4.1--P4.2, Bảng P4.2, P4.6}
\end{frame}

% SLIDE 27 -----------------------------------------------------------------
\begin{frame}{SETAR: dự báo test, khoảng và hạn chế}
\begin{columns}[T,onlytextwidth]
\begin{column}{.67\linewidth}
\bigfig{p4_3_setar_pi}{\linewidth}{4.02cm}
\centering{\scriptsize\color{midgray}Hình P4.3 -- Dự báo SETAR và khoảng kinh nghiệm 95\%}
\end{column}
\begin{column}{.30\linewidth}
\topic{Kết quả test}
\small\begin{tabular}{@{}lr@{}}\toprule
MAE & 19,4838\\
RMSE & 30,8688\\
RW RMSE & 30,0534\\
Bao phủ & 96,00\%\\
Độ rộng & 148,61\\\bottomrule
\end{tabular}
\vspace{0.18cm}\small Sai số test còn tự tương quan và hiệu ứng ARCH.
\end{column}
\end{columns}
\takeaway{Lợi thế trên validation không duy trì trên test: SETAR có cả MAE và RMSE cao hơn Random Walk.}
\src{Phần II, mục 7.4.3--7.4.6, Bảng P4.4--P4.10}
\end{frame}

\section{VI. Mô hình Markov Switching}
\partdivider{VI}{MÔ HÌNH\\MARKOV SWITCHING}{Trạng thái ẩn và chuyển đổi chế độ}{}

% MS-AR source: slides_ms_ar_usdvnd_2026.pdf (8 slides), also report part II 7.5.
% SLIDE M5-01 -------------------------------------------------------------
\begin{frame}{Markov Switching: mục tiêu và phạm vi}
\vspace{0.18cm}
\begin{columns}[T,onlytextwidth]
\begin{column}{.54\linewidth}
\topic{Bài toán nghiên cứu}
\small Mô tả sự thay đổi \textbf{chế độ biến động} của lợi suất log USD/VND; đồng thời lập dự báo tỷ giá \textbf{một phiên trước} trên tập test năm 2026.\par
\vspace{0.28cm}
\topic{Dữ liệu dùng để đánh giá}
\small Mẫu ước lượng đến hết năm 2025 (gộp train và validation); năm 2026 giữ làm \textbf{test ngoài mẫu}.
\end{column}
\begin{column}{.43\linewidth}
\begin{beamercolorbox}[wd=\linewidth,sep=.23cm,rounded=true]{block body}
\centering{\Huge\bfseries\color{darkred}MS--AR(2)}\\[.12cm]
{\small Hai trạng thái ẩn, tự hồi quy bậc hai}
\end{beamercolorbox}
\vspace{.26cm}
\small \textbf{$K=2$} được ấn định trước ước lượng; trạng thái biến động cao/thấp \textbf{không có nhãn sẵn} trong dữ liệu.
\end{column}
\end{columns}
\vspace{.32cm}
\takeaway{Đối chiếu khả năng dự báo một bước, nhưng vẫn phân biệt mục tiêu phân loại chế độ với độ chính xác mức tỷ giá.}
\src{Phần II, mục 7.5.1--7.5.3; MS--AR 2026, tr. 1}
\end{frame}

% SLIDE M5-02 -------------------------------------------------------------
\begin{frame}{Hai trạng thái ẩn và xác suất chuyển}
\begin{columns}[T,onlytextwidth]
\begin{column}{.52\linewidth}
\topic{Thiết kế mô hình}
\begin{itemize}\setlength{\itemsep}{4pt}
\item $S_t\in\{H,L\}$ không được quan sát trực tiếp.
\item Chọn \textbf{$K=2$} trước khi ước lượng; gắn nhãn cao/thấp \emph{sau ước lượng}, dựa vào $\sigma$.
\item Bộ lọc và bộ làm trơn xác định xác suất trạng thái theo tập thông tin khác nhau.
\end{itemize}
\vspace{.20cm}
\begin{tikzpicture}[node distance=1.18cm]
\node[flow,text width=1.63cm,minimum height=.74cm](h){\textbf{High}\\\scriptsize Biến động cao};
\node[flow,text width=1.63cm,minimum height=.74cm,right=of h](l){\textbf{Low}\\\scriptsize Biến động thấp};
\draw[arrow,bend left=25](h.north east) to node[above]{\scriptsize$11{,}06\%$}(l.north west);
\draw[arrow,bend left=25](l.south west) to node[below]{\scriptsize$8{,}63\%$}(h.south east);
\end{tikzpicture}
\end{column}
\begin{column}{.45\linewidth}
\topic{Ma trận chuyển tiếp ước lượng}
\[
\boxed{\mathbf P=\begin{pmatrix}
0{,}889435&0{,}110565\\
0{,}086334&0{,}913666
\end{pmatrix}}
\]
\vspace{.08cm}
{\scriptsize Hàng: $S_{t-1}$; cột: $S_t$ (thứ tự Cao, Thấp).}\\[.22cm]
\small\begin{tabular}{@{}lcc@{}}\toprule
 & Cao & Thấp\\\midrule
Xác suất duy trì & 88,94\% & 91,37\%\\
Kỳ vọng số phiên & 9,04 & 11,58\\\bottomrule
\end{tabular}
\end{column}
\end{columns}
\src{Phần II, mục 7.5.2--7.5.4, Bảng P5.2; MS--AR 2026, tr. 2}
\end{frame}

% SLIDE M5-03 -------------------------------------------------------------
\begin{frame}{Ma trận chuyển tiếp: cách đọc kết quả}
\centering
\[\boxed{\mathbf P=\begin{pmatrix}
0{,}889435&0{,}110565\\
0{,}086334&0{,}913666
\end{pmatrix}}\]
\vspace{.08cm}
\begin{columns}[T,onlytextwidth]
\begin{column}{.48\linewidth}
\begin{block}{Từ chế độ biến động cao}
\small\begin{itemize}\setlength{\itemsep}{5pt}
\item Giữ chế độ \textbf{cao}: $88{,}94\%$.
\item Chuyển sang \textbf{thấp}: $11{,}06\%$.
\item Thời gian kỳ vọng: \textbf{9,04 phiên}.
\end{itemize}
\end{block}
\end{column}
\begin{column}{.48\linewidth}
\begin{block}{Từ chế độ biến động thấp}
\small\begin{itemize}\setlength{\itemsep}{5pt}
\item Giữ chế độ \textbf{thấp}: $91{,}37\%$.
\item Chuyển sang \textbf{cao}: $8{,}63\%$.
\item Thời gian kỳ vọng: \textbf{11,58 phiên}.
\end{itemize}
\end{block}
\end{column}
\end{columns}
\vspace{.22cm}
\[
\mathbb{E}(D_i)=\frac{1}{1-p_{ii}}
\]
\notegray{Hai hàng tương ứng chế độ ở $t-1$, hai cột tương ứng chế độ ở $t$.}
\src{Phần II, mục 7.5.4, Bảng P5.2; MS--AR 2026, tr. 3}
\end{frame}

% SLIDE M5-04 -------------------------------------------------------------
\begin{frame}{MS--AR(2): đặc tả và cách dự báo}
\topic{Trung bình và phương sai thay đổi theo trạng thái}
\vspace{-.10cm}
\begin{align*}
r_t & = \mu_{S_t}+\phi_{1,S_t}\bigl(r_{t-1}-\mu_{S_{t-1}}\bigr) \\
    & \quad +\phi_{2,S_t}\bigl(r_{t-2}-\mu_{S_{t-2}}\bigr)+\varepsilon_t,\\[-1pt]
\varepsilon_t\mid S_t & \sim\mathcal N\bigl(0,\sigma_{S_t}^{2}\bigr).
\end{align*}
\vspace{-.08cm}
\begin{columns}[T,onlytextwidth]
\begin{column}{.48\linewidth}
\topic{Ước lượng tham số}
\small\begin{itemize}\setlength{\itemsep}{3pt}
\item $\mu$, AR(1), AR(2), $\sigma^2$ và ma trận $\mathbf P$ cùng ước lượng bằng \textbf{maximum likelihood}.
\item Xác suất \textbf{lọc}: chỉ sử dụng dữ liệu quan sát đến $t$.
\end{itemize}
\end{column}
\begin{column}{.48\linewidth}
\topic{Dự báo một bước}
\small\begin{itemize}\setlength{\itemsep}{3pt}
\item Trung bình điều kiện phụ thuộc $S_t,S_{t-1},S_{t-2}$.
\item Dùng xác suất trạng thái tại thời điểm dự báo; không dùng \textbf{xác suất làm trơn} hồi cứu.
\end{itemize}
\end{column}
\end{columns}
\vspace{.11cm}
\[\boldsymbol\pi_{t\mid t-1}=\boldsymbol\pi_{t-1\mid t-1}\mathbf P\]
\src{Phần I, mục 4.5; Phần II, mục 7.5.3; MS--AR 2026, tr. 4}
\end{frame}

% SLIDE M5-05 -------------------------------------------------------------
\begin{frame}{MS--AR(2): tham số ước lượng theo chế độ}
\vspace{.04cm}
\begin{columns}[T,onlytextwidth]
\begin{column}{.57\linewidth}
\topic{Trung bình và hệ số AR}
\begin{tikzpicture}[x=1cm,y=1cm,font=\scriptsize]
% axis mapping x=2.2 + 13.2*value, values between -0.06 to +0.13
\draw[midgray!60,line width=.4pt] (2.00,.12)--(2.00,2.82);
\foreach \v/\yy/\name in {0/2.4/$\mu$,1/1.52/AR(1),2/.64/AR(2)} {
  \node[anchor=east] at (.68,\yy) {\name};
}
% regime-high and regime-low bars; use scale 10 per unit, origin x=2.00
\fill[mainred] (2.00,2.54) rectangle (2.05244,2.77);
\fill[midgray!75] (1.98953,2.27) rectangle (2.00,2.49);
\fill[mainred] (2.00,1.66) rectangle (3.19936,1.89);
\fill[midgray!75] (1.46381,1.39) rectangle (2.00,1.61);
\fill[mainred] (1.54028,.78) rectangle (2.00,1.01);
\fill[midgray!75] (1.87513,.51) rectangle (2.00,.73);
\draw[midgray!28] (1.28,.2)--(1.28,2.82);
\draw[midgray!28] (2.72,.2)--(2.72,2.82);
\node[anchor=north] at (1.28,.2) {$-0{,}07$};
\node[anchor=north] at (2.00,.2) {$0$};
\node[anchor=north] at (2.72,.2) {$0{,}07$};
\node[anchor=west,fill=mainred,minimum width=.18cm,minimum height=.12cm,inner sep=0pt] at (.10,3.22) {};
\node[anchor=west] at (.42,3.22) {Cao};
\node[anchor=west,fill=midgray!75,minimum width=.18cm,minimum height=.12cm,inner sep=0pt] at (1.18,3.22) {};
\node[anchor=west] at (1.48,3.22) {Thấp};
\end{tikzpicture}
\end{column}
\begin{column}{.40\linewidth}
\topic{Độ lệch chuẩn theo chế độ}
\begin{tikzpicture}[x=1cm,y=1cm]
\node[anchor=west,font=\scriptsize] at (0,2.66) {Biến động cao};
\fill[mainred] (0,2.16) rectangle (2.52,2.49);
\node[anchor=west,font=\scriptsize] at (0,1.63) {Biến động thấp};
\fill[midgray!70] (0,1.14) rectangle (.421,1.47);
\node[anchor=west,font=\scriptsize] at (0,.71) {$\sigma_H\!\approx 5{,}99\,\sigma_L$};
\end{tikzpicture}
\end{column}
\end{columns}
\vspace{.04cm}
{\small\centering\renewcommand{\arraystretch}{1.15}
\begin{tabular}{@{}lrrrr@{}}\toprule
\textbf{Chế độ} & $\mu$ & AR(1) & AR(2) & $\sigma$\\\midrule
Cao & 0,005244 & 0,119936 & $-0{,}045972$ & 0,203432\\
Thấp & $-0{,}001047$ & $-0{,}053619$ & $-0{,}012487$ & 0,033954\\\bottomrule
\end{tabular}\par}
\vspace{.12cm}
\notegray{Tham số lợi suất và độ lệch chuẩn theo đơn vị điểm phần trăm; ước lượng đến 31/12/2025.}
\src{Phần II, mục 7.5.4, Bảng P5.1; MS--AR 2026, tr. 5}
\end{frame}

% SLIDE M5-06 -------------------------------------------------------------
\begin{frame}{MS--AR(2): khác biệt chủ yếu ở biến động}
\begin{columns}[T,onlytextwidth]
\begin{column}{.65\linewidth}
\bigfig{p5_1_regime}{\linewidth}{4.35cm}\par\vspace{0.03cm}\centering{\scriptsize\color{midgray}Hình P5.1 -- Lợi suất log và xác suất trạng thái trên mẫu ước lượng}\par
\end{column}
\begin{column}{.33\linewidth}
\topic{Đặc trưng chế độ}
\small\begin{tabular}{@{}lrr@{}}\toprule
 & Cao & Thấp\\\midrule
$\mu$ & 0,0052 & $-0{,}0010$\\
AR(1) & 0,1199 & $-0{,}0536$\\
AR(2) & $-0{,}0460$ & $-0{,}0125$\\
$\sigma$ & 0,2034 & 0,0340\\\bottomrule
\end{tabular}\par
\vspace{.14cm}
\footnotesize Độ phân tán là khác biệt nổi bật: $\sigma_H/\sigma_L\approx5{,}99$.
\end{column}
\end{columns}
\vspace{.06cm}
\takeaway{Xác suất làm trơn mô tả các cụm lịch sử rõ hơn; chỉ xác suất đã lọc / dự báo trước mới dùng cho dự báo tại thời điểm thực tế.}
\src{Phần II, mục 7.5.4, Hình P5.1; MS--AR 2026, tr. 6}
\end{frame}

% SLIDE M5-07 -------------------------------------------------------------
\begin{frame}{Markov: kiểm tra ngoài mẫu năm 2026}
\begin{columns}[T,onlytextwidth]
\begin{column}{.62\linewidth}
\bigfig{p5_3_markov_pi}{\linewidth}{4.43cm}\par\vspace{0.03cm}\centering{\scriptsize\color{midgray}Hình P5.3 -- Giá thực tế, dự báo một bước và khoảng 95\%}\par
\end{column}
\begin{column}{.36\linewidth}
\topic{175 phiên test 2026}
{\footnotesize\renewcommand{\arraystretch}{1.15}
\begin{tabular}{@{}lr@{}}\toprule
\textbf{Chỉ số} & \textbf{Kết quả}\\\midrule
MAE & 18,66\\
RMSE & 29,93\\
MAPE & 0,0713\%\\
Đúng chiều & 48,57\%\\
Bao phủ 95\% & 97,71\%\\
Độ rộng TB & 135,68\\\bottomrule
\end{tabular}\par}
\vspace{.10cm}
{\scriptsize\color{midgray}MAE, RMSE, độ rộng: VND/USD.}\par
\end{column}
\end{columns}
\vspace{.04cm}
\takeaway{Dự báo bám giá ở nhiều phiên nhưng dễ sai hơn trong pha biến động nhanh; tỷ lệ đúng chiều 48,57\% vẫn hạn chế.}
\src{Phần II, mục 7.5.6, Bảng P5.4; MS--AR 2026, tr. 7}
\end{frame}

% SLIDE M5-08 -------------------------------------------------------------
\begin{frame}{Markov Switching: kết luận và giới hạn}
\vspace{.08cm}
\begin{columns}[T,onlytextwidth]
\begin{column}{.475\linewidth}
\begin{block}{Kết quả đạt được}
\small\begin{itemize}\setlength{\itemsep}{4pt}
\item Phân biệt hai trạng thái chủ yếu theo độ biến động: $\sigma_H/\sigma_L\approx5{,}99$.
\item Trên test 2026, MAE $\approx18{,}66$ và RMSE $\approx29{,}93$ VND/USD.
\item Khoảng 95\% bao phủ $97{,}71\%$ quan sát.
\end{itemize}
\end{block}
\end{column}
\begin{column}{.475\linewidth}
\begin{block}{Giới hạn cần nêu rõ}
\small\begin{itemize}\setlength{\itemsep}{4pt}
\item Đúng chiều chỉ $48{,}57\%$; không có ưu thế rõ về hướng.
\item Phần dư ước lượng còn tự tương quan, không chuẩn và có ARCH.
\item Chưa có kiểm định tổn thất \emph{ghép cặp} với Random Walk.
\end{itemize}
\end{block}
\end{column}
\end{columns}
\vspace{.12cm}
\takeaway{RMSE thấp hơn Random Walk về số học, nhưng chưa đủ bằng chứng thống kê để kết luận MS--AR(2) vượt mô hình đối chứng.}
\src{Phần II, mục 7.5.5--7.5.6 và 8.4--8.7; MS--AR 2026, tr. 8}
\end{frame}

\section{VII. So sánh và kết luận}
\partdivider{VII}{SO SÁNH VÀ KẾT LUẬN}{Đánh giá dự báo, giới hạn và hướng hoàn thiện}{}

% SLIDE 32 -----------------------------------------------------------------
\begin{frame}{So sánh dự báo điểm trên 175 phiên test}
\vspace{0.14cm}
{\small\renewcommand{\arraystretch}{1.33}
\begin{tabularx}{\linewidth}{@{}X>{\raggedleft\arraybackslash}p{2.75cm}>{\raggedleft\arraybackslash}p{2.75cm}@{}}
\rowcolor{softred}\textbf{Mô hình} & \textbf{MAE} & \textbf{RMSE}\\\midrule
\textbf{Markov Switching AR(2)} & $18{,}6568$ & \textbf{$\approx29{,}93$}\\
ARMA--GARCH Student--$t$ & $18{,}5630$ & $30{,}0521$\\
\rowcolor{lightgray}\textbf{Random Walk (đối chứng)} & \textbf{$18{,}5629$} & $30{,}0534$\\
Kalman Local Level & $18{,}5678$ & $30{,}0545$\\
ARIMA(3,1,1) & $18{,}7363$ & $30{,}4392$\\
SETAR(6;4) & $19{,}4838$ & $30{,}8688$\\\bottomrule
\end{tabularx}}
\vspace{0.22cm}
\begin{columns}[T,onlytextwidth]
\begin{column}{.48\linewidth}
\begin{beamercolorbox}[wd=\linewidth,sep=.14cm,rounded=true]{block body}\small \textbf{MAE thấp nhất:} Random Walk.\end{beamercolorbox}
\end{column}
\begin{column}{.48\linewidth}
\begin{beamercolorbox}[wd=\linewidth,sep=.14cm,rounded=true]{block body}\small \textbf{RMSE thấp nhất:} Markov (chỉ số mô tả).\end{beamercolorbox}
\end{column}
\end{columns}
\src{Phần II, mục 8.4, Bảng 8.2; đơn vị sai số: VND/USD}
\end{frame}

% SLIDE 33 -----------------------------------------------------------------
\begin{frame}{So sánh khoảng dự báo và bằng chứng thống kê}
\vspace{0.02cm}
{\scriptsize\renewcommand{\arraystretch}{1.17}
\begin{tabularx}{\linewidth}{@{}X>{\raggedleft\arraybackslash}p{2.3cm}>{\raggedleft\arraybackslash}p{2.25cm}p{2.55cm}@{}}
\rowcolor{softred}\textbf{Mô hình} & \textbf{Bao phủ 95\%} & \textbf{Độ rộng TB} & \textbf{Loss ghép cặp với RW}\\\midrule
Random Walk & 94,86\% & 141,86 & Đối chứng\\
ARIMA & 95,43\% & 139,21 & Đã có HAC / bootstrap\\
ARMA--GARCH & 93,714\% & \textbf{94,53} & \textbf{Chưa có}\\
Kalman & 94,86\% & 141,57 & Đã có HAC / bootstrap\\
SETAR & 96,00\% & 148,61 & Đã có HAC / bootstrap\\
Markov & 97,71\% & 135,68 & \textbf{Chưa có}\\\bottomrule
\end{tabularx}}
\vspace{0.07cm}
\footnotesize\begin{itemize}\setlength{\itemsep}{1pt}
\item Khoảng GARCH hẹp nhất nhưng có tỷ lệ bao phủ thấp hơn danh nghĩa.
\item Kalman và RW có coverage và interval score gần nhau.
\end{itemize}
\takeaway{Độ rộng phải được đọc cùng coverage; kiểm định chưa đầy đủ không cho phép xếp hạng chắc chắn.}
\src{Phần II, mục 8.5--8.7, Bảng 8.1 và 8.4}
\end{frame}

% SLIDE 34 -----------------------------------------------------------------
\begin{frame}{Kết luận, giới hạn và hướng hoàn thiện}
\vspace{0.25cm}
\begin{columns}[T,onlytextwidth]
\begin{column}{.32\linewidth}
\begin{block}{Kết quả đạt được}
\small Năm mô hình mô tả các cấu trúc khác nhau: phụ thuộc tuyến tính, phương sai thay đổi, trạng thái ẩn và ngưỡng. Đã thống nhất tập test 175 phiên.
\end{block}
\end{column}
\begin{column}{.32\linewidth}
\begin{block}{Giới hạn bằng chứng}
\small Chưa chứng minh được lợi thế ổn định so với RW. Dữ liệu có 569 dòng OHLC bất nhất; một số kiểm định và tính tái lập chưa đầy đủ.
\end{block}
\end{column}
\begin{column}{.32\linewidth}
\begin{block}{Hướng hoàn thiện}
\small Xác minh nguồn Close; thống nhất mã/phiên bản và phép đánh giá; kiểm định ghép cặp còn thiếu; thử giai đoạn đánh giá mới.
\end{block}
\end{column}
\end{columns}
\vspace{0.42cm}
\takeaway{Chưa có bằng chứng nhất quán rằng các mô hình phức tạp cải thiện ổn định dự báo một bước so với Random Walk.}
\src{Phần II, mục 9.1--9.5}
\end{frame}

% Dedicated closing page matching the supplied reference.
\presentationThankYou

\end{document}
